%========================================================================
% CAPÍTULO 1 — INTRODUÇÃO
% Incluído por \input em main.tex. NÃO adicionar preâmbulo nem \begin{document}.
% Citações: só o conjunto canônico definido em sections/referencias.tex.
%========================================================================
\section{Introdução}\label{cap:intro}

A automação residencial deixou de ser um artigo de luxo e tornou-se um produto de
prateleira: no Brasil, uma lâmpada ou uma tomada Wi-Fi genérica --- em geral
construídas sobre a plataforma \textit{Tuya} --- custam entre R\$~30 e R\$~60 no
varejo. O mercado nacional de \textit{smart home} foi estimado em cerca de
US\$~2,7~bilhões em 2024, com projeção de crescimento próximo a 10\%~ao~ano até 2033
(IMARC, 2025), e a expansão brasileira, da ordem de 30\%~ao~ano, supera com folga a
média global de aproximadamente 12\%~ao~ano (IDC, 2024). O barateamento do hardware,
portanto, já ocorreu. Ainda assim, transformar esse hardware barato em automação
efetivamente utilizável permanece fora do alcance do domicílio típico brasileiro.

Os indicadores de adoção evidenciam a distância entre a disponibilidade do produto e
seu uso real. A proporção de domicílios brasileiros com ao menos um dispositivo
inteligente passou de 14,3\% em 2022 para 16,4\% em 2024 (IBGE, 2024), patamar muito
inferior aos cerca de 44\% observados nos Estados Unidos e no Reino Unido e à média
mundial de aproximadamente 33\% (STATISTA, 2023). A adoção nacional é, além de baixa,
profundamente desigual: 17,1\% nos domicílios urbanos contra 7,5\% nos rurais, e
19,8\% na região Sul contra 11,2\% no Nordeste (IBGE, 2024). A
Tabela~\ref{tab:penetracao} sintetiza esses recortes.

\begin{table}[H]
\centering
\caption{Penetração de dispositivos residenciais inteligentes: Brasil (recortes) e
referências internacionais.}
\label{tab:penetracao}
\begin{tabular}{@{}lcl@{}}
\toprule
\textbf{Recorte} & \textbf{Penetração (\%)} & \textbf{Fonte} \\
\midrule
Brasil --- domicílios (2022)     & 14,3        & IBGE, 2024 \\
Brasil --- domicílios (2024)     & 16,4        & IBGE, 2024 \\
Brasil --- área urbana           & 17,1        & IBGE, 2024 \\
Brasil --- área rural            & 7,5         & IBGE, 2024 \\
Brasil --- região Sul            & 19,8        & IBGE, 2024 \\
Brasil --- região Nordeste       & 11,2        & IBGE, 2024 \\
Estados Unidos / Reino Unido     & $\approx$44 & STATISTA, 2023 \\
Média mundial                    & $\approx$33 & STATISTA, 2023 \\
\bottomrule
\end{tabular}
\end{table}

Essa desigualdade não se explica por preço de compra, e sim por barreiras que se
acumulam depois da aquisição. A primeira é a conectividade: apenas 22\% dos
domicílios brasileiros dispõem de uma conexão à internet considerada satisfatória,
disparidade que se aprofunda por renda --- 73\% na classe A contra 3\% nas classes
D e~E (CETIC.br, 2024). Como a maioria dos dispositivos comerciais depende de
servidores do fabricante para funcionar, a casa ``inteligente'' deixa de operar
justamente onde a conectividade é mais precária. A segunda é demográfica: o Brasil
envelhece --- a parcela da população com 65 anos ou mais alcançou 10,9\% em 2022
(IBGE, 2024) --- e os idosos, que estão entre os que mais se beneficiariam de
interfaces por voz, são também os menos familiarizados com o letramento digital que a
instalação desses produtos exige. A terceira é a privacidade: o modelo comercial
dominante transmite continuamente áudio e hábitos domésticos a terceiros, o que
tensiona os princípios de minimização e finalidade da Lei Geral de Proteção de Dados
--- LGPD (BRASIL, 2018).

Convém, desde já, delimitar o que \textit{não} é o problema deste trabalho. O custo de
compra do dispositivo é tratado aqui como \textbf{premissa}, e não como a dor a ser
resolvida: o hardware Wi-Fi genérico já é acessível. O obstáculo decisivo à
democratização é outro --- é o \textbf{custo de comissionamento} (a instalação, o
pareamento e a configuração de cada aparelho, que exigem aplicativos proprietários,
criação de contas em nuvem e, para qualquer forma de controle local, o uso de portais
de desenvolvedor em inglês e ferramentas de linha de comando) somado à
\textbf{dependência estrutural de nuvem} para a operação cotidiana. É essa combinação,
e não o preço, que separa o hardware barato da automação utilizável. O princípio
operacional que orienta a proposta resume essa inversão de prioridades: ``nuvem uma
vez, local para sempre''.

\subsection{Problema de pesquisa}

O obstáculo central à democratização da automação residencial no Brasil não é o preço
do hardware --- já acessível ---, mas a \textbf{complexidade de instalar, conectar e
operar} os dispositivos, agravada pela dependência de nuvem e por aplicativos
proprietários. Cada dispositivo, ainda que barato, chega ao usuário atrelado ao
aplicativo e à conta em nuvem de seu fabricante; o controle local, quando existe,
depende de credenciais obtidas em portais de desenvolvedor voltados a especialistas.
O resultado é uma barreira de letramento digital que exclui precisamente os grupos que
mais se beneficiariam da tecnologia --- domicílios de menor renda, populações rurais e
do Nordeste, e a população idosa em crescimento. Formula-se, assim, o problema de
pesquisa: \textit{como tornar a automação residencial acessível a um usuário não
técnico, preservando a privacidade e o funcionamento sem dependência permanente de
nuvem?}

\subsection{Pergunta norteadora e hipótese}

Do problema deriva a seguinte \textbf{pergunta norteadora}: é viável um sistema de
automação residencial funcional, controlável por voz em português brasileiro,
\textit{local-first} e baseado em hardware de baixo custo, que possa ser
\textit{comissionado} por um usuário não técnico sem o aplicativo do fabricante e sem
portais de desenvolvedor?

Diferentemente de uma alegação socioeconômica ampla e não verificável, adota-se aqui
uma \textbf{hipótese de viabilidade falsificável dispositivo a dispositivo}, formulada
de modo a ser confirmada ou refutada por evidência empírica de cada aparelho:

\begin{quote}
\textit{É tecnicamente viável comissionar e controlar localmente, sem aplicativo do
fabricante e sem portais de desenvolvedor, dispositivos Wi-Fi genéricos de baixo
custo, com comandos de voz em português brasileiro processados integralmente no
\textit{hub}, a um custo incremental total inferior a R\$~200.}
\end{quote}

A hipótese é testada em dois dispositivos representativos do varejo brasileiro, com
status distintos: a tomada TP-Link Tapo~P110 --- na qual a forma plena foi
\textbf{confirmada}, com controle local demonstrado pelo protocolo KLAP a latências de
201 a 332~ms usando credenciais de consumidor (não de desenvolvedor) --- e a lâmpada
Intelbras/Tuya EWS~410 --- cujo comissionamento local permanece \textbf{em aberto (em
validação)}. Registre-se, desde a introdução, que a própria dificuldade de comissionar
a lâmpada Tuya, longe de enfraquecer o argumento, constitui \textbf{evidência empírica
da barreira} que este trabalho investiga: se a etapa se mostra custosa até para o
autor, a fortiori o é para o usuário não técnico. Os detalhes dessa evidência são
apresentados no Capítulo~6.

\subsection{Objetivos}

\subsubsection{Objetivo geral}

Demonstrar a viabilidade técnica de um \textit{hub} de automação residencial
\textit{local-first}, de baixo custo e controlável por voz em português brasileiro,
cujo mecanismo de comissionamento reduza a barreira de acesso do usuário não técnico.

\subsubsection{Objetivos específicos}

\begin{enumerate}
  \item Projetar uma arquitetura de software que isole a heterogeneidade de protocolos
        de dispositivos (Tuya, Tapo e outros) atrás de uma interface única, de modo a
        permitir a extensão do sistema sem reescrita do núcleo.
  \item Implementar o reconhecimento de fala em português brasileiro no próprio
        \textit{hub}, sem envio de áudio à nuvem, transcrevendo o comando localmente e
        descartando o áudio após a interpretação.
  \item Validar experimentalmente o controle local de hardware real, medindo latência
        (por percentis) e confiabilidade da comunicação com os dispositivos.
  \item Investigar e caracterizar um mecanismo de comissionamento de dispositivos
        Wi-Fi genéricos que dispense o aplicativo do fabricante e os portais de
        desenvolvedor.
  \item Definir e aplicar um protocolo de avaliação do artefato, contemplando custo
        total (\textit{Bill of Materials} --- BOM), acurácia de voz (taxa de erro de
        palavra --- WER --- e acurácia de intenção), consumo de energia e usabilidade
        junto ao público-alvo (\textit{System Usability Scale} --- SUS).
        % [...] A coletar: BOM datada em 2 cenários (hub reaproveitado × completo);
        % WER sobre corpus pt-BR (>=50 enunciados, >=3 falantes); SUS com n=3-5.
\end{enumerate}

\subsection{Justificativa}

A relevância deste trabalho decorre da convergência entre uma necessidade social e uma
lacuna técnica. Do lado social, os dados apresentados evidenciam uma adoção baixa e
desigual (IBGE, 2024; STATISTA, 2023), condicionada por conectividade precária em boa
parte dos domicílios (CETIC.br, 2024) e por uma transição demográfica que amplia a
população idosa (IBGE, 2024) --- público para o qual a interação por voz representa uma
porta de entrada natural, desde que não exija letramento técnico para a instalação.
Some-se a isso a tensão entre o modelo comercial dependente de nuvem e os princípios da
LGPD (BRASIL, 2018), que torna desejável um arranjo em que os dados domésticos não
saiam da residência.

Do lado técnico, existem soluções maduras de controle local --- notadamente o
\textit{Home Assistant} (HOME ASSISTANT, 2025) --- e um estado da arte de
comissionamento padronizado no protocolo Matter (CONNECTIVITY STANDARDS ALLIANCE,
2024). Contudo, essas soluções ou pressupõem hardware de rádio adicional e ainda caro
ou escasso no varejo popular brasileiro (o caso do Matter), ou, no caso do
\textit{Home Assistant} aplicado ao Wi-Fi barato, pressupõem o \textbf{pareamento
prévio pelo aplicativo do fabricante} --- exatamente a etapa que exclui o usuário não
técnico. A contribuição do DOMUS não está em disputar a sofisticação linguística dos
assistentes comerciais, e sim em atacar o \textbf{comissionamento acessível e local} de
dispositivos Wi-Fi genéricos, alinhado ao paradigma \textit{local-first} descrito por
KLEPPMANN et al. (2019), no qual o usuário é o dono de seus dados e a nuvem é
conveniência, não requisito. A viabilidade do processamento de voz no próprio
\textit{hub}, sem GPU, apoia-se em modelos de reconhecimento de fala robustos como o
\textit{Whisper} (RADFORD et al., 2022), que operam em hardware modesto. O trabalho
adota, ainda, as diretrizes da ABNT para trabalhos acadêmicos (ABNT, 2011) na sua
estruturação e apresentação.

\subsection{Organização do trabalho}

Além desta introdução, o trabalho está organizado em mais seis capítulos. O
\textbf{Capítulo~2 --- Fundamentação teórica} revisa os conceitos de automação
residencial e IoT, a fragmentação de protocolos, o paradigma \textit{local-first}, o
reconhecimento de fala local, a acessibilidade e a segurança e privacidade sob a LGPD.
O \textbf{Capítulo~3 --- Trabalhos relacionados} confronta o DOMUS com os assistentes
comerciais dependentes de nuvem, com o \textit{Home Assistant} e com o protocolo
Matter, delimitando o recorte da contribuição. O \textbf{Capítulo~4 --- Materiais e
métodos} descreve a abordagem metodológica (ciência de projeto e desenvolvimento
orientado a especificações), o padrão \textit{Adapter} como mecanismo contra a
fragmentação, as tecnologias empregadas e os critérios e instrumentos de avaliação. O
\textbf{Capítulo~5 --- Desenvolvimento} detalha a arquitetura e a implementação do
\textit{hub}, do \textit{pipeline} de voz e do mecanismo de comissionamento. O
\textbf{Capítulo~6 --- Resultados e discussão} apresenta as medições obtidas
--- controle local do Tapo~P110, \textit{pipeline} de voz e o estado dos critérios de
sucesso --- e discute seus limites. Por fim, o \textbf{Capítulo~7 --- Considerações
finais} retoma a tese, expõe as limitações e aponta os trabalhos futuros.
